\chapter{Implications and Limitations}
\label{chap:discussion}

\section{Implications for using Expanded Welch}

Normal Wald remains the simpler baseline; the results do not support
Expanded Welch as a general replacement. Expanded Welch is useful where
its larger cutoff brings excessive false positives closer to 5\% at
an acceptable cost in power. Where Wald already rejects too rarely,
the same change makes rejection more conservative. The comparison
should also report valid-result frequency, because failed calculations
can lower the rejection rate without providing a test decision.

For an observed feature--label comparison, report the estimated MI
difference, its standard error and the p-value together. The difference
gives the estimated size and direction of the change; the standard error
gives the sampling uncertainty assigned by the method. The sparse-table
results show why the direction matters: unequal estimation biases can
shift or reverse the estimated difference, as
Equation~\eqref{eq:focused-null-bias-shift} illustrates. If Expanded
Welch is used, report its estimated degrees of freedom as well, since
they determine the rejection cutoff for that comparison.

Before applying either test, assess its performance on populations that
match the application's category definitions, sample sizes and margin
patterns. Include plausible population MI values and distinct equal-MI
tables, then compare false positives, power and usable-result frequency.
This assessment is particularly important for sparse data, since neither
method is reliable in every studied setting. Choose the method using
that evidence before examining the application's observed p-values.

\begin{samepage}
\section{Limits of the evidence}

The rejection rates and Monte Carlo intervals apply to the constructed
population pairs. The intervals quantify uncertainty from the finite number
of simulated samples; they do not describe how performance varies across
untested populations.
\par
\end{samepage}

The primary populations use a copying mixture with identical row and
column margins within each population. This gives a clear way to vary
association without changing category frequencies. The four-cell check
supports the representative \(4\times4\) method comparison and
reproduces the main \(16\times16\) estimation failures at its matched
settings, although power changes numerically with the construction.
Copying favours matching categories in both
populations, whereas the four-cell transfer concentrates association
in the same chosen block. The different-skew regime changes category
frequencies while keeping those association locations aligned. Its
results therefore demonstrate that distinct equal-MI tables can be
difficult to compare, but do not describe every way two tables with
equal MI can differ. A broader assessment would vary association
locations, population MI values and margins, including unequal
row and column margins within a population.

Equal sample sizes simplify the main studies by making the two leading
bias corrections cancel in \(\Deltahat\). The remaining difference
bias in the experiment consequently comes from estimation errors not
removed by the leading MI bias correction. With unequal sample sizes, the
corrections no longer cancel, and the two contributions to the standard
error and degrees of freedom can become less balanced. The derivation
allows unequal sample sizes. The supporting fixed-total check assesses
both allocation directions in two \(4\times4\) different-skew pairs.
There, imbalance changes bias, uncertainty and power, while the
larger Expanded Welch cutoff lowers already-conservative null rates.
This extends the comparison beyond equal samples, but a wider allocation
study would be needed for other table sizes and population structures.

The sampling assumptions also limit transfer to other applications.
Paired observations would require a covariance term in the variance of
the difference, whereas the present samples are independent. Continuous
MI estimators have different estimation errors from fixed categorical
tables. Similarly, increasing the number of categories with sample size
changes the fixed-alphabet setting of the derivation. The empirical
recommendation concerns the independent categorical comparisons studied
here, at a two-sided 5\% threshold. The online-shopping application
illustrates interpretation on observed data, but its unknown population
MI values do not allow a direct measurement of false-positive rates or power.

Exact independence needs a different mathematical approximation because
the first-order MI variance vanishes there. The population-MI
study also shows that positive population MI is not enough to make that
approximation accurate at a particular sample size: small positive MI
can remain close to the problematic region. This explains why improving
the rejection cutoff and improving MI's sampling approximation near independence
are separate tasks.
Section~\ref{sec:independence-boundary} derives the second-order
approximation needed there.
